\section{Factorization of Polynomials over a Field}

\exercise{23.1}{Polynomial division 1}
Work in $\Z_{7}[x]$, reducing coefficients after each subtraction. The inverse of the divisor's leading coefficient is $1$. Repeatedly subtract the multiple that cancels the current highest power. The complete long-division steps are
\[
\begin{array}{c|l}
\text{term added to quotient}&\text{new remainder}\\\hline
x^4&x^5+3x^4+4x^2+4x+2\\
x^3&x^4+3x^3+4x^2+4x+2\\
x^2&x^3+4x+2\\
x&5x^2+2\\
5&4x+3\\
\end{array}
\]
Thus $\boxed{q(x)=x^4+x^3+x^2+x+5}$ and $\boxed{r(x)=4x+3}$. The final remainder has degree below the divisor's degree; reversing the subtractions verifies $f=gq+r$.

\exercise{23.2}{Polynomial division 2}
Work in $\Z_{7}[x]$, reducing coefficients after each subtraction. The inverse of the divisor's leading coefficient is $5$. Repeatedly subtract the multiple that cancels the current highest power. The complete long-division steps are
\[
\begin{array}{c|l}
\text{term added to quotient}&\text{new remainder}\\\hline
5x^4&x^4+4x^2+4x+2\\
5x^2&4x^3+5x^2+4x+2\\
6x&x+2\\
\end{array}
\]
Thus $\boxed{q(x)=5x^4+5x^2+6x}$ and $\boxed{r(x)=x+2}$. The final remainder has degree below the divisor's degree; reversing the subtractions verifies $f=gq+r$.

\exercise{23.3}{Polynomial division 3}
Work in $\Z_{11}[x]$, reducing coefficients after each subtraction. The inverse of the divisor's leading coefficient is $6$. Repeatedly subtract the multiple that cancels the current highest power. The complete long-division steps are
\[
\begin{array}{c|l}
\text{term added to quotient}&\text{new remainder}\\\hline
6x^4&3x^4+3x+6\\
7x^3&4x^3+3x+6\\
2x^2&9x^2+3x+6\\
10x&4x+6\\
2&4\\
\end{array}
\]
Thus $\boxed{q(x)=6x^4+7x^3+2x^2+10x+2}$ and $\boxed{r(x)=4}$. The final remainder has degree below the divisor's degree; reversing the subtractions verifies $f=gq+r$.

\exercise{23.4}{Polynomial division 4}
Work in $\Z_{11}[x]$, reducing coefficients after each subtraction. The inverse of the divisor's leading coefficient is $9$. Repeatedly subtract the multiple that cancels the current highest power. The complete long-division steps are
\[
\begin{array}{c|l}
\text{term added to quotient}&\text{new remainder}\\\hline
9x^2&3x^3+x^2\\
5x&6x^2+x\\
10&2\\
\end{array}
\]
Thus $\boxed{q(x)=9x^2+5x+10}$ and $\boxed{r(x)=2}$. The final remainder has degree below the divisor's degree; reversing the subtractions verifies $f=gq+r$.

\exercise{23.5}{All multiplicative generators modulo 5}
The group has order $4$. The element $2$ has this full order: it satisfies $2^{4}=1$, and the checks $2^{2}\equiv4\pmod{5}$ exclude every proper divisor of the order (check exponents $(4)/q$ for each prime $q$ dividing $4$). In a cyclic group of order $4$, the generators are exactly $g^k$ with $\gcd(k,4)=1$. Here those exponents are $1,3$, yielding the complete generator set
\[\boxed{\{2,3\}}.
\]

\exercise{23.6}{All multiplicative generators modulo 7}
The group has order $6$. The element $3$ has this full order: it satisfies $3^{6}=1$, and the checks $3^{3}\equiv6, 3^{2}\equiv2\pmod{7}$ exclude every proper divisor of the order (check exponents $(6)/q$ for each prime $q$ dividing $6$). In a cyclic group of order $6$, the generators are exactly $g^k$ with $\gcd(k,6)=1$. Here those exponents are $1,5$, yielding the complete generator set
\[\boxed{\{3,5\}}.
\]

\exercise{23.7}{All multiplicative generators modulo 17}
The group has order $16$. The element $3$ has this full order: it satisfies $3^{16}=1$, and the checks $3^{8}\equiv16\pmod{17}$ exclude every proper divisor of the order (check exponents $(16)/q$ for each prime $q$ dividing $16$). In a cyclic group of order $16$, the generators are exactly $g^k$ with $\gcd(k,16)=1$. Here those exponents are $1,3,5,7,9,11,13,15$, yielding the complete generator set
\[\boxed{\{3,5,6,7,10,11,12,14\}}.
\]

\exercise{23.8}{All multiplicative generators modulo 23}
The group has order $22$. The element $5$ has this full order: it satisfies $5^{22}=1$, and the checks $5^{11}\equiv22, 5^{2}\equiv2\pmod{23}$ exclude every proper divisor of the order (check exponents $(22)/q$ for each prime $q$ dividing $22$). In a cyclic group of order $22$, the generators are exactly $g^k$ with $\gcd(k,22)=1$. Here those exponents are $1,3,5,7,9,13,15,17,19,21$, yielding the complete generator set
\[\boxed{\{5,7,10,11,14,15,17,19,20,21\}}.
\]

\exercise{23.9}{Four linear factors modulo five}
Since $4=-1$ in $\Z_5$,
\[x^4+4=x^4-1=(x^2-1)(x^2+1).
\]
Now $x^2-1=(x-1)(x+1)$ and $x^2+1=x^2-4=(x-2)(x+2)$. Therefore the requested linear factorization is
\[\boxed{(x-1)(x+1)(x-2)(x+2)}.
\]

\exercise{23.10}{Three linear factors modulo seven}
Substitution shows that $x=-1$ is a zero. Division gives
\[x^3+2x^2+2x+1=(x+1)(x^2+x+1).
\]
In $\Z_7$, $(x-2)(x+3)=x^2+x-6=x^2+x+1$. Thus
\[\boxed{x^3+2x^2+2x+1=(x+1)(x-2)(x+3)}.
\]
Expansion checks all coefficients modulo $7$.

\exercise{23.11}{A nonmonic cubic modulo eleven}
The residues $3,4,-3$ are zeros, suggesting the three factors. Their product expands as
\[2(x-4)(x-3)(x+3)=2(x-4)(x^2-9)
=2x^3-8x^2-18x+72.
\]
Reducing modulo $11$ gives $2x^3+3x^2-7x-5$, as required. Hence this is the complete linear factorization, including its leading coefficient $2$.

\exercise{23.12}{A reducible cubic modulo five}
The value at $x=2$ is $8+4+3=15=0$ modulo $5$, so the cubic is not irreducible. Division by $x-2$ gives $x^2+2x+1=(x+1)^2$. Therefore
\[\boxed{x^3+2x+3=(x-2)(x+1)^2}\quad\text{in }\Z_5[x].
\]
All three linear factors are irreducible, with a repeated factor explicitly recorded.

\exercise{23.13}{An irreducible cubic modulo five}
A cubic over a field is reducible exactly when it has a zero in the field. The values of $2x^3+x^2+2x+2$ at $0,1,2,3,4$ are $2,2,1,1,4$ in $\Z_5$, none zero. It is therefore irreducible. If a monic irreducible factor is desired, write it as
\[\boxed{2(x^3+3x^2+x+1)}.
\]
The scalar $2$ is a unit, and the cubic factor is irreducible; there is no factorization into positive lower degrees.

\exercise{23.14}{A quadratic over three fields}
The discriminant of $x^2+8x-2$ is $8^2-4(1)(-2)=72$. Its roots are $-4\pm3\sqrt2$, neither rational. A quadratic over a field is reducible if and only if it has a root there, so it is irreducible over $\Q$. Both roots are real, however, giving
\[x^2+8x-2=(x+4-3\sqrt2)(x+4+3\sqrt2).
\]
Thus it is reducible over $\R$ and consequently over $\C$.

\exercise{23.15}{A quadratic with negative discriminant}
For $x^2+6x+12$, the discriminant is $36-48=-12$. There is no real root, hence no rational root. It is irreducible over both $\Q$ and $\R$. Over $\C$, its roots are $-3\pm i\sqrt3$, so
\[x^2+6x+12=(x+3-i\sqrt3)(x+3+i\sqrt3)
\]
is a factorization and it is reducible there.

\exercise{23.16}{A cubic over the rationals}
A reducible cubic over $\Q$ would have a rational root. The rational-root theorem for the monic polynomial $x^3+3x^2-8$ limits candidates to $\pm1,\pm2,\pm4,\pm8$. Their values, in that order, are
\[-4,-6,12,-4,104,-24,696,-328,
\]
all nonzero. Thus there is no rational root and the cubic is irreducible over $\Q$.

\exercise{23.17}{A quartic with no quadratic factor}
The only possible rational roots of the monic integer polynomial are $\pm1$, and both give $-20$, so no linear factor exists. By Gauss's lemma, any remaining rational factorization could be taken as
\[(x^2+ax+b)(x^2+cx+d),\qquad a,b,c,d\in\Z.
\]
Comparing the $x^3$ and constant coefficients gives $c=-a$ and $bd=1$, so $(b,d)=(1,1)$ or $(-1,-1)$. The $x^2$ coefficient is $b+d-a^2=-22$. In the first case this forces $a^2=24$; in the second it forces $a^2=20$. Neither is an integer square. Thus there is no quadratic factorization either, and $x^4-22x^2+1$ is irreducible over $\Q$.

\exercise{23.18}{An Eisenstein quadratic}
For $x^2-12$, use the prime $p=3$. It does not divide the leading coefficient $1$, divides all other coefficients (including the zero linear coefficient), and $3^2=9$ does not divide $-12$. Eisenstein's criterion therefore applies and proves irreducibility over $\Q$.

\exercise{23.19}{An Eisenstein cubic}
Use $p=3$: it does not divide the leading coefficient $8$, it divides the other coefficients $6,-9,24$, and $9$ does not divide the constant coefficient $24$. Thus all Eisenstein hypotheses hold and the cubic is irreducible over $\Q$.

\exercise{23.20}{An Eisenstein polynomial of degree ten}
Use $p=3$ for $4x^{10}-9x^3+24x-18$. Although $3\nmid4$ and $3$ divides all lower coefficients, $3^2=9$ divides the constant coefficient $-18$. This prime fails the square condition. No other prime divides all the lower nonzero coefficients, since their gcd is $3$. Thus the polynomial satisfies no Eisenstein criterion in its displayed form. Failure of the criterion alone does not decide reducibility.

\exercise{23.21}{Another Eisenstein polynomial of degree ten}
For $2x^{10}-25x^3+10x^2-30$, the prime $5$ divides every lower coefficient but not the leading coefficient $2$. Its square $25$ does not divide $-30$. Thus Eisenstein applies with $p=5$, proving irreducibility over $\Q$.

\exercise{23.22}{All rational zeros of a quartic}
The rational-root theorem permits only fractions with numerator dividing $10$ and denominator dividing $6$. Testing candidates finds $x=-5/2$ and $x=2/3$. Exact division, or multiplication, gives
\[6x^4+17x^3+7x^2+x-10=(2x+5)(3x-2)(x^2+x+1).
\]
The final quadratic has discriminant $1-4=-3$, so it has no rational zero. Consequently the complete rational zero set is $\boxed{\{-5/2,\ 2/3\}}$. The displayed factorization checks the candidate search and excludes further rational roots.

\exercise{23.23}{Correcting irreducibility}
A nonconstant polynomial $f\in F[x]$ is irreducible over $F$ if it cannot be written as $f=gh$ with \emph{both} $g,h\in F[x]$ of positive degree. The printed definition forbids even $f=1\cdot f$, which every polynomial has. Trivial factorizations by nonzero constant units must be allowed.

\exercise{23.24}{An equivalent irreducibility definition}
The stated definition is correct. A nonconstant $f$ is nonzero. In any factorization $f=gh$, both factors are nonzero, so a factor belonging to $F$ is a nonzero constant and hence a unit in $F[x]$. Requiring such a factor in every factorization is equivalent to forbidding a product of two positive-degree polynomials.

\exercise{23.25}{True or false: irreducibility and root bounds}
\begin{enumerate}[label=(\alph*)]
\item True: a degree-one polynomial cannot be a product of two positive-degree polynomials.
\item True for the same reason; the nonzero scalar factor $3$ is a unit in $\Q[x]$.
\item True: $\sqrt3$ is irrational, so this quadratic has no rational root.
\item False: in $\Z_7$, $2^2+3=7=0$, so $x^2+3$ is reducible.
\item True: degree additivity forces any polynomial unit over a field to be a nonzero constant.
\item True: the source repeats the same assertion as in (e), so the same reasoning applies.
\item True: a polynomial of degree $n$ is nonzero and has at most $n$ distinct zeros over a field.
\item True: regard it as the same nonzero degree-$n$ polynomial over the extension field $E$ and apply the root bound there.
\item True: $ax+b$, $a\ne0$, has zero $-b/a\in F$.
\item False: over an infinite field such as $\Q$, the zero polynomial vanishes at infinitely many inputs. The assertion becomes true if one requires the polynomial to be nonzero.
\end{enumerate}

\exercise{23.26}{Primes for which a specified linear factor divides}
The factor theorem says that $x+2=x-(-2)$ divides the polynomial modulo $p$ exactly when its value at $-2$ is zero modulo $p$. This value is
\[16-8+4+2+1=15.
\]
Thus $p$ must divide $15$. The prime possibilities are exactly $\boxed{3,5}$, and both do satisfy the condition.

\exercise{23.27}{Irreducible quadratics over two elements}
Every monic quadratic has form $x^2+ax+b$. It has no root at $0$ only if $b=1$; then its value at $1$ is $a$, which is nonzero only if $a=1$. Thus the only irreducible quadratic is $\boxed{x^2+x+1}$. There is no further nonzero scalar multiple to list because $1$ is the only nonzero scalar in $\Z_2$.

\exercise{23.28}{Irreducible cubics over two elements}
A cubic is irreducible exactly when it has no field root. Write it $x^3+ax^2+bx+c$. Avoiding root $0$ requires $c=1$; avoiding root $1$ requires $a+b\ne0$, so $(a,b)=(1,0)$ or $(0,1)$. The complete list is
\[\boxed{x^3+x^2+1,\qquad x^3+x+1}.
\]

\exercise{23.29}{Irreducible quadratics over three elements}
Test the monic quadratics $x^2+ax+b$ at $0,1,2$. The monic ones with no zero are exactly
\[x^2+1,\quad x^2+x+2,\quad x^2+2x+2.
\]
Their respective value triples are $(1,2,2)$, $(2,1,2)$, and $(2,2,1)$. All quadratic leading coefficients are $1$ or $2$, and multiplying by a nonzero constant preserves irreducibility. Thus the full answer is these three polynomials and twice each of them, six in total.

\exercise{23.30}{Irreducible cubics over three elements}
A cubic is irreducible exactly when none of $0,1,2$ is a zero. The monic list, together with values at these inputs, is
\[
\begin{array}{l|ccc}h(x)&h(0)&h(1)&h(2)\\\hline
x^3+2x^2+1&1&1&2\\
x^3+2x^2+x+1&1&2&1\\
x^3+2x+1&1&1&1\\
x^3+x^2+2x+1&1&2&2\\
x^3+x^2+2&2&1&2\\
x^3+x^2+x+2&2&2&1\\
x^3+2x+2&2&2&2\\
x^3+2x^2+2x+2&2&1&1
\end{array}.
\]
To see completeness, write a monic cubic as $x^3+ax^2+bx+c$. Necessarily $c=1$ or $2$; testing the nine pairs $(a,b)$ for each $c$ and discarding zeros at $1$ or $2$ gives precisely these eight rows. The full list of irreducible cubics consists of these eight polynomials and their multiples by $2$, sixteen in all.

\exercise{23.31}{Counting irreducible quadratics over a prime field}
There are $p^2$ monic quadratics. A reducible monic quadratic is uniquely $(x-a)(x-b)$ for an unordered pair of field elements, allowing $a=b$. There are $\binom p2+p=p(p+1)/2$ such pairs. Therefore the number of monic irreducible quadratics is
\[p^2-\frac{p(p+1)}2=\frac{p(p-1)}2.
\]
Each has $p-1$ distinct nonzero scalar multiples, and every nonmonic quadratic is a unique scalar times a monic one. Thus the requested count of \emph{all} irreducible quadratics is $\boxed{p(p-1)^2/2}$.

\exercise{23.32}{Synopsis of the bound on roots}
Each distinct zero supplies a distinct linear factor by the factor theorem, and repeatedly removing these factors lowers the degree by one each time; a nonzero degree-$n$ polynomial cannot supply more than $n$ such factors.

\exercise{23.33}{Synopsis of cyclicity of finite field subgroups}
By the finite abelian group theorem, write the finite multiplicative subgroup with invariant factors $d_1\mid\cdots\mid d_r$. Every element then satisfies $x^{d_r}=1$. The polynomial root bound gives $|G|\le d_r$, while the product of invariant factors is $|G|\ge d_r$. Equality forces all earlier factors to be trivial, so the group is cyclic.

\exercise{23.34}{A reducible prime-power polynomial}
In $\Z_p$, Fermat gives $a^p=a$, and all intermediate binomial coefficients of exponent $p$ vanish. Hence
\[(x+a)^p=x^p+a^p=x^p+a.
\]
Since $p\ge2$, this is a product of $p$ linear factors, so it is not irreducible. This includes $a=0$, when the polynomial is $x^p$.

\exercise{23.35}{The reciprocal polynomial}
Let $g(x)=a_n+a_{n-1}x+\cdots+a_0x^n$. Since $a\ne0$, substitute $a^{-1}$ and factor out $a^{-n}$:
\[g(a^{-1})=a^{-n}(a_0+a_1a+\cdots+a_na^n)=a^{-n}f(a)=0.
\]
Thus $1/a$ is a zero of the reciprocal polynomial, with the nonzero assumption justifying all inverse powers.

\exercise{23.36}{The remainder theorem}
Division by $x-a$ gives $f(x)=(x-a)q(x)+r(x)$, where $r=0$ or $\deg r<1$. Thus $r$ is a constant. Evaluating at $a$ yields $f(a)=0\cdot q(a)+r=r$. Hence the remainder is exactly the constant polynomial $\boxed{f(a)}$.

\exercise{23.37}{Reduction of integer coefficients}
\textbf{(a)} Reduce each coefficient modulo $m$. Sums are preserved coefficientwise; products are preserved because each coefficient of a product is a finite sum of coefficient products, and integer reduction preserves both operations. Every polynomial over $\Z_m$ can be lifted by choosing integer representatives for its finitely many coefficients, so the map is onto.\par
\textbf{(b)} Suppose instead that $f$ is reducible over $\Q$. By Gauss's lemma it has a factorization $f=gh$ in $\Z[x]$ with $0<\deg g,\deg h<n$. The leading coefficient of $f$ remains nonzero modulo $m$ by the degree hypothesis. Because it is the product of the leading coefficients of $g,h$, neither of those can become zero modulo $m$. Their reduced degrees are therefore unchanged. Reducing $f=gh$ then gives the forbidden factorization into two degrees below $n$, a contradiction. Thus $f$ is irreducible over $\Q$.\par
\textbf{(c)} Modulo $5$, the polynomial becomes $x^3+2x+1$. Its values at $0,1,2,3,4$ are $1,4,3,4,3$, none zero. As a cubic over the field $\Z_5$, it is irreducible, and its degree remains $3$. Part (b) proves that $x^3+17x+36$ is irreducible over $\Q$.
